% !TeX root = expose.tex
% !TeX spellcheck = en_US

\subsection*{Situation and Problem Statement}
In many virtual reality (VR) games and applications, immersion, the feeling of truly being present and engaged in the virtual world, plays a major role. While modern VR headsets excel at providing high-fidelity visual and auditory stimuli, they often lack adequate tactile feedback. This absence can quickly break the illusion and destroy immersion when a player tries to touch a virtual object but feels nothing but empty air. Because active, motorized haptic controllers are often too complex, expensive, or restrictive to implement seamlessly, there is a need to explore alternative ways to trick the player's brain into perceiving physical interaction.

\subsection*{Project Focus and Planned Methodology}
To address this missing physical feedback, this project focuses on the concept of "passive haptics" and "visuo-haptic illusions." This approach uses inert, real-world objects as physical proxies for virtual counterparts, allowing players to physically touch what they visually perceive. The core question is whether integrating these real-world objects successfully tricks players into a higher state of immersion and makes them believe these objects belong to the virtual experience.

We will develop a modular VR prototype encompassing various experimental setups. While the exact scenarios will be determined during the course of the project, they will involve pairing everyday physical objects with distinct virtual experiences to measure the impact on the user's suspension of disbelief. Another question we aim to explore is whether player emotion (e.g., fear, disgust, etc.) can be influenced by the presence of passive haptics.

To evaluate the efficacy of these setups, we will conduct user tests comparing the passive haptics condition against a standard baseline. The evaluation will rely on standardized HCI immersion metrics, such as the Igroup Presence Questionnaire (IPQ), which is widely used to explicitly measure a user's perception of spatial presence, involvement, and the overall "realness" of the VR environment \cite{presence-questionnaires}. To measure emotional arousal and valence, we will use the Self-Assessment Manikin (SAM) \cite{sam}, a questionnaire that allows participants to rate their emotional response on dimensions of arousal and valence. Additionally, we will measure physiological responses, such as heart rate variability (HRV) to provide objective data on the user's emotional state during the VR experience.

\subsection*{Comparison to Related Work}
This project builds upon and distinguishes itself from several foundational studies in the field. First, we draw on the 2015 study "Substitutional Reality" by Simeone et al. \cite{substitutional-reality}, which established that a physical proxy does not need to be a perfect 1:1 replica of the virtual object (e.g., using a physical umbrella as a virtual lightsaber) to maintain user engagement. In their research, they mapped physical objects in the real environment to different virtual settings, such as a spaceship or a medieval courtyard, demonstrating how everyday items can be repurposed to interact with the virtual world.

Additionally, we will integrate insights and methodologies from recent studies focusing on emotional and physiological responses in VR. For example, a study by Kobayashi et al. \cite{Kobayashi.2025} compared identically designed real and virtual rooms to evaluate spatial perception using methods including questionnaires and HRV. The results demonstrated that while real rooms were perceived as more pleasant and preferred, the VR counterparts triggered higher physiological arousal. This highlights that VR provokes strong arousal and cognitive processing effort, making it not emotionally equivalent to reality.

Similarly, a study by Khundrakpam et al. \cite{Khundrakpam.2025} focused on a VR platform designed to systematically generate and measure stress responses using methods including PPG/HRV and mood scales. Through interactions like hitting colored cylinders or jogging on a VR treadmill, the study proved that VR exercises are highly engaging and effectively trigger intended stress demands, with HRV providing valuable insights into these stress reactions.

Furthermore, the study by Marocco et al. \cite{Marocco.2025} investigated whether virtual nature spaces provide emotional regeneration after stressful VR scenarios. Utilizing questionnaires and HRV measurements, the findings revealed an interesting nuance: higher physiological activation in natural VR scenes was actually evaluated positively by users, demonstrating that lower physiological activation does not automatically equate to a better emotional effect.

Finally, Zhang et al. \cite{Zhang.2025} examined in their study how selection of colors in a virtual workspace serves as an emotional stimulus. Using measurements including HRV and IPQ, researchers found that red environments generated anxiety and high arousal, yellow generated happiness, and blue induced a calm, lower-arousal state which demonstrates that not only virtual content of scenes can influence emotions.

By adapting these measurement approaches and understanding the complex physiological arousal inherent to VR, our project will assess whether tangible physical proxies can further manipulate, amplify, or ground these intense emotional states.

\subsection*{Expected Results}
The expected outcome is a functioning prototype demonstrating the selected interactive scenarios, alongside a comprehensive empirical evaluation of the user experience. By systematically gathering subjective and behavioral data, the project aims to scientifically demonstrate whether integrating everyday physical objects successfully heightens the player's immersion, emotional arousal, and sense of presence compared to traditional VR setups.